%!TEX root = ../Thesis.tex

\chapter{Stand der Technik und Wissenschaft}\label{cap:wissensstand}
\markboth{Stand der Technik und Wissenschaft}{Stand der Technik und Wissenschaft}

\improvement[inline]{Kapitel 2 strukturieren und für jeden Abschnitt einen zu erkedigenden Punkt einfügen - arbeit/chapter2.tex}

\section{Business Intelligence}\label{sec:bi} 

\subsection{Begriffe, Ziele und Nutzen}
Unter \glsentrylong{bi} (\gls{bi}) wird ein Oberbegriff verstanden, der Anwendungen, Infrastrukturen und Werkzeuge sowie bewährte Vorgehensweisen umfasst, um den Zugriff auf Informationen und deren Analyse zur Verbesserung und Optimierung von Entscheidungen und Leistungsfähigkeit zu ermöglichen \cite{GartnerABI}. \gls{bi} wird damit als Kombination aus technischen Komponenten und organisatorischer Nutzungsperspektive beschrieben, weil neben Werkzeugen auch Praktiken der Informationsnutzung explizit Bestandteil der Definition sind \cite{GartnerABI}. 
Ein integriertes Verständnis von \gls{bi} und Analytics adressiert die IT-basierte Entscheidungsunterstützung als Gesamtkonzept, das fachliche, technische und organisatorische Aspekte gemeinsam berücksichtigt \cite{BaarsKemperSpringerBook}. Dieses Gesamtkonzept schließt etablierte Ansätze aus dem Data-Warehouse- und Reporting-Umfeld ein und berücksichtigt darüber hinaus weiterführende Analytics-Ansätze \cite{BaarsKemperSpringerBook}. 
\gls{bi} kann als durchgängiger Prozess beschrieben werden, der von der Datenbereitstellung über die Datenanalyse bis zur Unternehmenssteuerung und zur Informationsverteilung reicht \cite{MuellerLenzMPDL}. Zu den in diesem Prozess verorteten Bausteinen zählen Data-Warehousing für die Datenbereitstellung sowie OLAP und Data Mining für die Analyse, wobei diese Elemente als integrale Bestandteile des BI-Ansatzes ausgewiesen werden \cite{MuellerLenzMPDL}. Die mit \gls{bi} adressierten Managementaufgaben umfassen unter anderem Prognose, Planung, Controlling und Monitoring, die im Kontext analytischer Informationssysteme als typische Anwendungsfelder der BI-gestützten Unternehmenssteuerung beschrieben werden \cite{MuellerLenzMPDL}. Der Nutzen von \gls{bi} wird in dieser Logik primär über die verbesserte Entscheidungs- und Steuerungsfähigkeit operationalisiert, weil die Analyse von Informationen explizit auf die Optimierung von Entscheidungen und Performance ausgerichtet ist \cite{GartnerABI}. 
Für \glspl{kmu} werden Datenanalytik und ihre Nutzung als leistungsrelevant eingeordnet, weil sie als potenzieller Beitrag zur Performance und zur Produktivitätssteigerung von \glspl{kmu} beschrieben werden \cite{OECDDataAnalyticsSMEs}. Gleichzeitig werden für \glspl{kmu} Hemmnisse der Nutzung von Datenanalytik adressiert, weil interne und externe Barrieren der Adoption und der Nutzung von Datenanalytik in \glspl{kmu} explizit diskutiert werden \cite{OECDDataAnalyticsSMEs}. Eine praxisnahe Einordnung für \glspl{kmu} beschreibt zudem, dass historische Daten zur fundierteren Entscheidungsfindung als relevant wahrgenommen werden, während fehlendes Know-how und begrenzte finanzielle Ressourcen die Umsetzung wirksamer BI-Lösungen erschweren können \cite{Raj2018SMECaseStudy}. 
Für die weitere Arbeit sind Begriffe, Ziele und Nutzen von \gls{bi} damit als Grundlage relevant, um BI-Maßnahmen konsistent bewerten und in Entscheidungslogiken zur Technologieintegration einordnen zu können.


\subsection{Architekturen und Datenbereitstellung}

\subsubsection{Data Warehouse}
Ein \gls{dwh} wird als Datenbasis verstanden, die Managemententscheidungen unterstützt und dafür Daten themenorientiert, integriert, zeitorientiert und nicht flüchtig bereitstellt \cite{Inmon1996}. Die Themenorientierung beschreibt, dass Daten entlang fachlicher Gegenstandsbereiche organisiert werden, um Analysen und Berichte entlang unternehmerischer Fragestellungen zu ermöglichen \cite{Inmon1996}. Die Integration adressiert die Vereinheitlichung heterogener Datenquellen, sodass Begriffe, Schlüssel und Werte domänenübergreifend konsistent interpretierbar sind \cite{OracleDWHConcepts}. Die Zeitorientierung bedeutet, dass Daten im \gls{dwh} historisiert werden und damit zeitbezogene Vergleiche und Trendanalysen unterstützen \cite{Inmon1996}. Die Nichtflüchtigkeit beschreibt, dass Daten im \gls{dwh} typischerweise nicht durch operative Transaktionen überschrieben werden, sondern stabil für Analysezwecke bereitgestellt werden \cite{OracleDWHConcepts}. 

Für die Modellierung analytischer Daten im \gls{dwh} wird häufig die dimensionale Modellierung genutzt, bei der Fakten über Kennzahlen mit Dimensionen für Kontextinformationen kombiniert werden \cite{KimballRoss2013}. Ein verbreitetes Strukturprinzip ist das Sternschema, in dem eine Faktentabelle über Schlüsselbeziehungen mit mehreren Dimensionstabellen verbunden wird \cite{KimballRoss2013}. Dieses Modellierungsprinzip zielt darauf ab, Abfragen für Reporting und Analyse zu vereinfachen und die Verständlichkeit für fachliche Nutzerinnen und Nutzer zu erhöhen \cite{KimballRoss2013}. Die Architekturperspektive eines \gls{dwh} umfasst damit sowohl die zentrale Datenhaltung als auch die semantische Strukturierung der Daten zur Nutzung in analytischen Anwendungen \cite{BaarsKemper2021}. 

\subsubsection{Data Lake}
Ein Data Lake wird als Konzept beschrieben, das aus einer Sammlung von Speichereinheiten für verschiedene Datenbestände besteht, die zusätzlich zu den ursprünglichen Quellsystemen vorgehalten werden \cite{GartnerDataLake}. Die in einem Data Lake abgelegten Daten werden in einer nahezu exakten oder exakten Kopie des Quellformats gespeichert, wodurch die ursprüngliche Struktur weitgehend erhalten bleibt \cite{GartnerDataLake}. Ein zentrales Merkmal ist, dass Daten zunächst in ihrem ursprünglichen, untransformierten Zustand gespeichert werden und Transformationen erst bei Bedarf erfolgen \cite{MSDataLakeGuide}. Dieses Vorgehen wird als \emph{schema-on-read} beschrieben, weil die Strukturierung und Interpretation beim Auslesen und nicht zwingend beim Laden festgelegt wird \cite{MSDataLakeGuide}. Demgegenüber wird für klassische \gls{dwh}-Ansätze ein \emph{schema-on-write} beschrieben, bei dem Struktur und Transformationen bereits beim Ingest in die Datenbasis angewandt werden \cite{MSDataLakeGuide}. 

Der Data Lake adressiert die Speicherung strukturierter, semi-strukturierter und unstrukturierter Daten in einem gemeinsamen Ablagekonzept, um eine breite Nutzung für Analysezwecke zu ermöglichen \cite{AWSDataLake}. Die funktionale Einordnung des Data Lakes als „Single Store“ für Rohdaten zur späteren Analyse wird als Bestandteil moderner Datenpipelines diskutiert \cite{Fowler2015DataLake}. Gleichzeitig wird hervorgehoben, dass der Nutzen eines Data Lakes von Governance, Definition, Sicherheit und Nutzungsorientierung abhängt, weil die reine Zusammenführung disparater Daten ohne Steuerungsmechanismen zu Fehlannahmen über den Wert der Datenhaltung führen kann \cite{GartnerDataLakeFallacy}. 

\subsubsection{ETL und ELT}
Die Datenbereitstellung für \gls{bi}-Systeme basiert typischerweise auf Datenintegrationsprozessen, die Daten aus Quellsystemen extrahieren, transformieren und in ein Zielsystem laden \cite{IBMETL}. \gls{etl} wird dabei als Prozess beschrieben, der Daten aus mehreren Quellen kombiniert, bereinigt und organisiert, um ein konsistentes Ziel-Dataset in einem \gls{dwh}, Data Lake oder einem vergleichbaren Zielsystem bereitzustellen \cite{IBMETL}. Die Transformationsschritte umfassen in diesem Verständnis insbesondere Harmonisierung, Qualitätsbereinigung und Strukturierung, die für Analyse- und Reportingzwecke erforderlich sind \cite{IBMETL}. Ein \gls{etl}-Vorgehen transformiert Daten vor dem Laden in die Zieldatenplattform, wodurch das Zielsystem bereits kuratierte und strukturierte Daten erhält \cite{DatabricksETLvsELT}. 

\gls{elt} wird als Variante beschrieben, bei der Daten zunächst geladen und anschließend innerhalb des Zielsystems transformiert werden \cite{DatabricksETLvsELT}. 
Dieses Vorgehen wird insbesondere mit leistungsfähigen, skalierbaren Zielplattformen begründet, weil Transformationen dann die Rechenkapazitäten des Zielsystems nutzen \cite{DatabricksETLvsELT}. Die Wahl zwischen \gls{etl} und \gls{elt} wird als Architekturentscheidung eingeordnet, die von Anforderungen an Datenlatenz, Governance, Transformationskomplexität und Zielplattform-Fähigkeiten abhängt \cite{BaarsKemper2021}. Für die Architektur von BI-Lösungen bedeutet dies, dass \gls{dwh}- und Data-Lake-Konzepte jeweils unterschiedliche Integrations- und Transformationslogiken nahelegen, die in der Datenbereitstellung konsistent umgesetzt werden müssen \cite{MSDataLakeGuide}. Vor diesem Hintergrund ist die klare Abgrenzung von Speicher- und Integrationsprinzipien eine Voraussetzung, um Datenbereitstellung so zu gestalten, dass Reporting und weiterführende Analytik auf einer nachvollziehbaren und wartbaren Datenbasis aufsetzen können \cite{BaarsKemper2021}. 


\subsection{Analytische Funktionen}

\subsubsection{Reporting}
Reporting bezeichnet die Fähigkeit, formatierte und interaktive Berichte zu erstellen und diese skalierbar zu verteilen sowie zu terminieren \cite{GartnerInformationDelivery}. 
Reporting wird als Bestandteil der Informationsbereitstellung in Analytics- und \gls{bi}-Plattformen eingeordnet, weil entsprechende Plattformen Berichte und Dashboards als typische Ausprägungen der Information Delivery vorsehen \cite{GartnerABIPlatformReq,GartnerInformationDelivery}. Business-Intelligence-Reporting wird als Nutzung von Visualisierungen beschrieben, um Ergebnisse leichter verständlich und teilbar zu machen \cite{MicrosoftWhatIsBI}. Als Reporting-Formen werden interaktive Dashboards, Diagramme, Graphen und Karten genannt, die Anwenderinnen und Anwendern eine aktuelle Sicht auf das Geschäftsgeschehen ermöglichen \cite{MicrosoftWhatIsBI}. Dashboards werden als Teilmenge des Reportings beschrieben, die Informationen in intuitiven Darstellungen bereitstellt und Zustände von Leistungskennzahlen im Verhältnis zu Zielwerten sichtbar macht \cite{GartnerInformationDelivery}. Dashboards werden zudem mit der Verbreitung von Echtzeitdaten aus operativen Anwendungen in Verbindung gebracht, wodurch Reporting nicht ausschließlich rückblickend angelegt ist \cite{GartnerInformationDelivery}. 

\subsubsection{Self-Service}
Self-Service Analytics wird als Form von \gls{bi} beschrieben, in der Fachbereiche befähigt und ermutigt werden, Abfragen auszuführen und Berichte eigenständig zu erzeugen, wobei nur begrenzte IT-Unterstützung erforderlich ist \cite{GartnerSelfServiceAnalytics}. Self-Service Analytics wird dabei häufig durch einfach zu bedienende Werkzeuge mit grundlegenden analytischen Fähigkeiten sowie ein für Verständlichkeit und Datenzugriff vereinfachtes Datenmodell charakterisiert \cite{GartnerSelfServiceAnalytics}. Ad-hoc Query wird als Fähigkeit beschrieben, die auch als Self-Service Reporting bezeichnet wird und Nutzerinnen und Nutzern ermöglicht, eigene Fragen an die Daten zu stellen, ohne dass die IT einen Bericht erstellen muss \cite{GartnerInformationDelivery}. Für diese Fähigkeit wird eine robuste semantische Schicht gefordert, damit Anwenderinnen und Anwender verfügbare Datenquellen navigieren können \cite{GartnerInformationDelivery}. 
Für die nachhaltige Nutzbarkeit werden zudem Query-Governance und Auditing als Anforderungen genannt, um eine performante Ausführung von Abfragen sicherzustellen \cite{GartnerInformationDelivery}. 
Analytics- und \gls{bi}-Plattformen werden in diesem Kontext als Systeme beschrieben, die nichttechnischen Nutzergruppen Self-Service-Zugriff auf freigegebene Daten sowie Analyse-, Visualisierungs- und Reporting-Funktionen bereitstellen \cite{GartnerABIPlatformReq}. 

\subsubsection{OLAP}
Online Analytical Processing (OLAP) wird als Tätigkeit beschrieben, bei der Text- und Zahlendaten aus Data Warehouses und Data Marts abgefragt und für analytische Zwecke präsentiert werden \cite{JukicJukicMalliaris2008OLAP}. OLAP wird als Technologie für schnelle, komplexe Abfragen beziehungsweise multidimensionale Analysen großer Datenmengen in einem Data Warehouse, einem Data Lake oder einem anderen Datenrepositorium beschrieben \cite{IBMOlap}. Die multidimensionale Perspektive wird mit der Zerlegung betrieblicher Daten in mehrere Dimensionen begründet, die beispielsweise Zeit, Standort oder Produktkategorien repräsentieren können \cite{IBMOlap}. OLAP wird als Bestandteil von \gls{bi} und Decision Support sowie als Grundlage für Forecasting- und Reporting-Anwendungsfälle eingeordnet \cite{IBMOlap}. Zur Unterstützung interaktiver Analysen wird im OLAP-Kontext die Pivot-Operation beschrieben, bei der Dimensionen zur Aggregation einer Kennzahl ausgewählt und aggregierte Werte in einer tabellarischen Sicht dargestellt werden \cite{ChaudhuriDayal1997}. Als weitere zentrale Operationen werden Roll-up und Drill-down beschrieben, wobei Roll-up eine zusätzliche Gruppierung entlang einer Dimension und Drill-down die Umkehrung dieser Verdichtung darstellt \cite{ChaudhuriDayal1997}. Slice-and-dice wird als Operation beschrieben, die die Dimensionalität reduziert, indem eine Projektion auf Teilmengen von Dimensionen für ausgewählte Werte anderer Dimensionen vorgenommen wird \cite{ChaudhuriDayal1997}. 

\subsubsection{Advanced Analytics}
Advanced Analytics wird als autonome oder teilautonome Untersuchung von Daten oder Inhalten mit fortgeschrittenen Techniken und Werkzeugen beschrieben, die typischerweise über traditionelle \gls{bi}-Ansätze hinausgehen \cite{GartnerAdvancedAnalytics}. Als Ziele werden die Entdeckung tieferer Erkenntnisse, die Erstellung von Vorhersagen und die Generierung von Empfehlungen genannt \cite{GartnerAdvancedAnalytics}. Als beispielhafte Techniken werden unter anderem Data- und Text-Mining, Machine Learning, Forecasting sowie multivariate Statistik, Simulation und komplexe Ereignisverarbeitung aufgeführt \cite{GartnerAdvancedAnalytics}. Advanced Analytics wird in Abgrenzung zu rein deskriptiver Analytik mit Fragestellungen verknüpft, die Gründe erklären, zukünftige Entwicklungen antizipieren und Handlungsmöglichkeiten ableiten \cite{GartnerPressReleaseAdvancedAnalytics2014}. In Plattformkontexten wird zudem die Fähigkeit beschrieben, automatisierte Insights mithilfe von Machine-Learning-Techniken für Endnutzerinnen und Endnutzer zu generieren \cite{GartnerABIPlatformReq}. 


\subsection{Datenqualität, Governance und Sicherheit}

\subsubsection{Datenqualität}
Datenqualität wird als Eignung von Daten für den vorgesehenen Verwendungszweck beschrieben \cite{WangStrong1996}. 
Diese Eignung ist kontextabhängig, weil Anforderungen an Datenqualität von der jeweiligen Aufgabe und den Erwartungen der Datennutzenden abhängen \cite{WangStrong1996}. Datenqualität wird als multidimensionales Konstrukt eingeordnet, weil unterschiedliche Qualitätsaspekte gleichzeitig relevant sein können \cite{WangStrong1996}. Wang und Strong strukturieren Datenqualitätsdimensionen in intrinsische, kontextuelle, repräsentationsbezogene und zugangsbezogene Kategorien \cite{WangStrong1996}. Die Dimensionen Genauigkeit, Vollständigkeit, Konsistenz und Aktualität werden als zentrale Datenqualitätsmerkmale in diesem Rahmen diskutiert \cite{WangStrong1996}. 

ISO/IEC~25012 definiert ein allgemeines Datenqualitätsmodell für Daten, die in strukturierter Form in einem Computersystem vorgehalten werden \cite{ISO25012}. Das Modell ist darauf ausgelegt, Anforderungen an Datenqualität zu etablieren, Datenqualitätsmaße zu definieren sowie Datenqualitätsevaluationen zu planen und durchzuführen \cite{ISO25012}. Für BI-Kontexte ist eine solche Operationalisierung relevant, weil analytische Auswertungen auf konsistenten und interpretierbaren Daten beruhen und Datenqualität als Voraussetzung für belastbare Informationsbereitstellung adressiert wird \cite{BaarsKemper2021}. 

\subsubsection{Data Governance}
Data Governance wird als Ausübung von Autorität und Kontrolle über das Management von Daten beschrieben \cite{KhatriBrown2010}. Khatri und Brown unterscheiden Governance und Management, indem Governance die zu treffenden Entscheidungen und die Zuordnung von Verantwortlichkeiten beschreibt, während Management die Umsetzung dieser Entscheidungen umfasst \cite{KhatriBrown2010}. Data Governance wird damit als Mechanismus eingeordnet, der Entscheidungsrechte, Verantwortlichkeiten und Leitplanken für datenbezogene Entscheidungen festlegt \cite{KhatriBrown2010}. Khatri und Brown leiten hierfür mehrere Datenentscheidungsdomänen ab und verknüpfen diese mit der Frage, wer innerhalb einer Organisation welche Datenentscheidungen verantwortet \cite{KhatriBrown2010}. Eine Governance-Perspektive adressiert damit explizit die organisatorische Verankerung von Datenverantwortung, die für konsistente Definitionen, kontrollierten Datenzugriff und die nachhaltige Nutzung gemeinsamer Datenbestände erforderlich ist \cite{KhatriBrown2010}. 

Für BI-Architekturen ist Data Governance zudem als Steuerungsrahmen relevant, weil BI- und Analytics-Fähigkeiten typischerweise fachbereichsübergreifende Datennutzung, einheitliche Kennzahldefinitionen und abgestimmte Datenbereitstellung erfordern \cite{BaarsKemper2021}. Die Ausrichtung von Datenbereitstellung und Datennutzung auf explizite Zuständigkeiten wird in diesem Kontext als Voraussetzung dafür eingeordnet, dass Daten konsistent interpretiert und wiederverwendbar bereitgestellt werden können \cite{KhatriBrown2010}. 

\subsubsection{Sicherheit und Datenschutz}
Informationssicherheit wird als Schutz von Informationen und Informationssystemen vor unbefugtem Zugriff, Nutzung, Offenlegung, Störung, Veränderung oder Zerstörung beschrieben \cite{NISTInfoSec}. Die Schutzziele Vertraulichkeit, Integrität und Verfügbarkeit werden als grundlegende Sicherheitsziele zur Einordnung dieser Schutzanforderungen verwendet \cite{NISTCIA}. Vertraulichkeit wird dabei als Eigenschaft definiert, die unbefugte Offenlegung von Informationen verhindert \cite{NISTCIA}. Integrität wird als Eigenschaft definiert, die Genauigkeit und Vollständigkeit von Informationen schützt \cite{NISTCIA}. Verfügbarkeit wird als Eigenschaft definiert, die den bedarfsgerechten Zugriff autorisierter Entitäten auf Informationen und Systeme sicherstellt \cite{NISTCIA}. 

Für BI-Systeme ist Informationssicherheit besonders relevant, weil BI-Architekturen Daten konsolidieren und damit Schutzbedarfe aus verschiedenen Quellsystemen in einer gemeinsamen Datenbasis zusammenführen \cite{BaarsKemper2021}. Ein risikobasierter Ansatz zur Steuerung von Informationssicherheitsmaßnahmen wird in ISO/IEC~27001 im Rahmen eines Information Security Management Systems (ISMS) beschrieben, das die Vertraulichkeit, Integrität und Verfügbarkeit von Informationen durch einen Risikomanagementprozess erhält \cite{ISO27001}. Technische und organisatorische Sicherheitsmaßnahmen werden in diesem Verständnis als Bestandteil eines Managementprozesses verstanden, der auf die Reduktion relevanter Risiken ausgerichtet ist \cite{ISO27001}. 

Datenschutz ist in BI-Kontexten relevant, weil BI-Anwendungen auch personenbezogene Daten verarbeiten können und deren Verarbeitung an rechtliche Anforderungen gebunden ist \cite{GDPR_EURLex}. Die DSGVO stellt hierfür einen zentralen europäischen Regulierungsrahmen dar, der Anforderungen an die Verarbeitung personenbezogener Daten sowie an deren Schutz adressiert \cite{GDPR_EURLex}. Für die BI-Gestaltung ergibt sich daraus die Notwendigkeit, Datenzugriffe, Verarbeitungszwecke und Schutzmaßnahmen so zu steuern, dass Informationssicherheitsziele und datenschutzrechtliche Anforderungen zusammenwirken \cite{ISO27001,GDPR_EURLex}. 

Datenqualität, Data Governance und Sicherheit bilden damit komplementäre Gestaltungsfelder, weil Datenqualität die Verlässlichkeit analytischer Ergebnisse beeinflusst, Governance die kontrollierte und nachvollziehbare Datennutzung strukturiert und Sicherheit den Schutz der verarbeiteten Informationen absichert \cite{WangStrong1996,KhatriBrown2010,ISO27001}. 


\subsection{Herausforderungen und Erfolgsfaktoren in \glspl{kmu}}

\subsubsection{Herausforderungen}
\gls{bi}-Initiativen in \glspl{kmu} sind häufig durch begrenzte finanzielle Ressourcen, begrenzte Zeitbudgets und knappe personelle Kapazitäten geprägt \cite{BianchiniMichalkova2019}. Finanzielle Restriktionen werden als zentrale Hürde für die Einführung von \gls{bi} in \glspl{kmu} beschrieben \cite{TaticDzaficHaracic2018}. 
Unzureichendes Wissen und fehlende Erfahrung im Umfeld von \gls{bi} werden als Barriere für die Einführung und Nutzung von \gls{bi}-Systemen in \glspl{kmu} benannt \cite{TaticDzaficHaracic2018}. Eine fehlende Vision und nicht klar definierte strategische Ziele werden als Hemmnis für die Implementierung von \gls{bi} in \glspl{kmu} aufgeführt \cite{TaticDzaficHaracic2018}. Unklar definierte KPIs werden als weitere Barriere genannt, weil sie die Ausrichtung von Analysen, Simulationen und Berichten erschweren \cite{TaticDzaficHaracic2018}. 

Herausforderungen werden zudem in der Bedienbarkeit und im Aufbau von Berichten verortet, wenn Systeme als zu komplex wahrgenommen werden \cite{GudfinnssonStrand2017}. 
Kompetenzdefizite und ein daraus abgeleiteter Trainingsbedarf werden als wiederkehrende Herausforderung für \gls{bi}-Adoption in \glspl{kmu} beschrieben \cite{GudfinnssonStrand2017}. Probleme in der Datenqualität werden als Hindernis genannt, wenn widersprüchliche Datenbestände und unzureichende Datenbasis die Nutzung analytischer Ergebnisse beeinträchtigen \cite{GudfinnssonStrand2017}. Schnittstellenprobleme zwischen Quellsystemen und analytischen Anwendungen werden als Herausforderung beschrieben, weil eingeschränkte Exportmöglichkeiten manuelle Importprozesse nach sich ziehen können \cite{GudfinnssonStrand2017}. Manuelle Import- und Exportschritte werden als Ursache für höhere Aufwände und für eine geringere Aktualität von Analysedaten eingeordnet \cite{GudfinnssonStrand2017}. 

\subsubsection{Erfolgsfaktoren}
Erfolgsfaktoren für \gls{bi} in \glspl{kmu} werden in organisatorische, prozessuale und technologische Kategorien strukturiert \cite{OlszakZiemba2012}. 
Top-Management-Unterstützung wird als relevanter Erfolgsfaktor beschrieben, weil fehlende Unterstützung als wesentliche organisatorische Barriere im SME-Kontext identifiziert wird \cite{OlszakZiemba2012}. Wissen über Fähigkeiten und Nutzen von \gls{bi}-Systemen wird als Erfolgsbedingung eingeordnet, weil fehlendes Wissen als zentrale Hürde für \gls{bi}-Vorhaben in \glspl{kmu} adressiert wird \cite{OlszakZiemba2012}. Klare fachliche Problemdefinitionen werden als Erfolgsfaktor beschrieben, weil unklare Business-Probleme als Hemmnis für die Zielausrichtung einer \gls{bi}-Initiative benannt werden \cite{OlszakZiemba2012}. Die Abstimmung zwischen fachlichen Anforderungen und BI-Systemausgestaltung wird als Erfolgsfaktor eingeordnet, weil fehlende Alignment-Aspekte als Barriere in \glspl{kmu} dokumentiert werden \cite{OlszakZiemba2012}. 

Benutzerorientierte Qualifizierung wird als Erfolgsfaktor beschrieben, weil fehlendes User-Training und fehlende Unterstützung als Barrieren in \glspl{kmu} genannt werden \cite{OlszakZiemba2012}. Widerstände von Anwenderinnen und Anwendern werden als Barriere beschrieben, wodurch Change- und Enablement-Maßnahmen als Erfolgsfaktoren für Akzeptanz ableitbar sind \cite{OlszakZiemba2012}. Budgetüberschreitungen und unzureichendes Projektmanagement werden als wiederkehrende Problembereiche beschrieben, wodurch realistische Planung und kontrollierte Umsetzung als Erfolgsdimensionen nahegelegt werden \cite{OlszakZiemba2012}. Externe Unterstützung wird als pragmatische Option eingeordnet, wenn internes Know-how begrenzt ist und die Organisation dennoch in der Lage sein muss, Implementierung und Betrieb bedarfsgerecht zu steuern \cite{TaticDzaficHaracic2018}. Eine konsistente Verankerung von Datenanalytik in \glspl{kmu} wird als leistungsrelevant eingeordnet, weil Datenanalytik mit Produktivitäts- und Performancepotenzialen für \glspl{kmu} verknüpft wird \cite{BianchiniMichalkova2019}. 

Für diese Arbeit sind die genannten Herausforderungen und Erfolgsfaktoren als Kontext relevant, weil sie Entscheidungskriterien für eine ressourcensensitive, strategische Einführung von \gls{bi} in \glspl{kmu} begründen.


\section{Prozessautomatisierung}\label{sec:procauto}

\subsection{Begriffe und Abgrenzungen}
Prozessautomatisierung wird im Kontext des \emph{Business Process Management} (BPM) als IT-gestützte Umsetzung verstanden, bei der Arbeit in Organisationen so gestaltet und unterstützt wird, dass konsistente Ergebnisse erzielt und Verbesserungspotenziale genutzt werden können \cite{DumasEtAl2013}. BPM wird als „art and science“ der Steuerung beschrieben, wie Arbeit in einer Organisation ausgeführt wird, und bezieht sich dabei auf ganze Ketten von Ereignissen, Aktivitäten und Entscheidungen, die Wert für Organisation und Kundschaft erzeugen \cite{DumasEtAl2013}. Der Begriff „Prozess“ wird in diesem Zusammenhang als Bezeichnung für solche Wertschöpfungsketten verwendet \cite{DumasEtAl2013}. Prozessautomatisierung adressiert damit die technische Unterstützung der Ausführung dieser Prozesse durch Informationssysteme \cite{DumasEtAl2013}. 

Der Begriff \emph{Workflow} wird als computergestützte Unterstützung oder Automatisierung eines Geschäftsprozesses im Ganzen oder in Teilen definiert, bei der Dokumente, Informationen oder Aufgaben nach prozeduralen Regeln zwischen Beteiligten zur Bearbeitung weitergegeben werden \cite{Aalst2003WorkflowTutorial}. Ein \emph{Workflow-Management-System} wird als System definiert, das Workflows mithilfe von Software definiert, erstellt und deren Ausführung steuert, wobei es Prozessdefinitionen interpretiert, mit Beteiligten interagiert und bei Bedarf IT-Werkzeuge oder Anwendungen aufruft \cite{Aalst2003WorkflowTutorial}. Die Workflow-Perspektive fokussiert damit die Ausführung und Orchestrierung von Arbeitsschritten entlang einer modellierten Ablaufstruktur \cite{Aalst2003WorkflowTutorial}. Ein Workflow-Management-System wird zudem als System beschrieben, das prozedurale Automatisierung eines Geschäftsprozesses bereitstellt, indem es die Abfolge von Arbeitsaktivitäten verwaltet und geeignete personelle oder IT-Ressourcen für die jeweiligen Schritte aktiviert \cite{OMGWfMF12}. Der Begriff \emph{Workflow-Automatisierung} kann in dieser Logik als Automatisierung von Prozessschritten verstanden werden, die durch ein Workflow-System kontrolliert und verteilt werden \cite{OMGWfMF12}. 

\emph{Robotic Process Automation} (RPA) wird als Sammelbegriff für Werkzeuge beschrieben, die über die Benutzeroberfläche anderer Informationssysteme in einer Weise operieren, wie es ein Mensch tun würde \cite{AalstBichlerHeinzl2018RPA}. RPA wird als Automatisierung in einem „outside-in“-Ansatz beschrieben, bei dem bestehende Informationssysteme unverändert bleiben und menschliche Interaktionen durch Softwareagenten ersetzt werden \cite{AalstBichlerHeinzl2018RPA}. RPA wird in Abgrenzung zu anderen Automatisierungstechnologien als „lightweight“ eingeordnet, weil der Zugriff typischerweise über die Präsentationsschicht erfolgt, um menschliches Verhalten nachzuahmen \cite{HermJanieschRausch2021}. RPA wird außerdem als Prozessautomatisierungstechnologie für hochvolumige, repetitive und regelbasierte Aufgaben beschrieben \cite{HindelCabreraStierle2020}. Diese Charakterisierung grenzt RPA von klassischer Workflow- beziehungsweise Backend-Automatisierung ab, weil RPA primär die Ausführung bestehender Benutzerinteraktionen automatisiert, während Workflow-Management und backendseitige Integration die Prozessausführung stärker über definierte Prozessmodelle, Systemlogik und Schnittstellen gestalten \cite{AalstBichlerHeinzl2018RPA,OMGWfMF12}. 

Für die weitere Arbeit sind diese Abgrenzungen relevant, weil sie eine konsistente Einordnung der betrachteten Automatisierungsansätze und ihrer Einsatzlogik ermöglichen.


\subsection{Prozessidentifikation und -auswahl}
Die Identifikation und Auswahl geeigneter Prozesse stellt im \emph{Business Process Management} (BPM) einen frühen und entscheidenden Schritt dar, weil Verbesserungs- und Automatisierungsmaßnahmen typischerweise auf konkrete, abgegrenzte Prozesssegmente angewendet werden \cite{DumasEtAl2013}. 
BPM beschreibt Prozesse als Ketten von Ereignissen, Aktivitäten und Entscheidungen, die in einer Organisation Wert erzeugen, wodurch die Auswahl eines „Prozesses“ stets auch eine Festlegung des Betrachtungsumfangs impliziert \cite{DumasEtAl2013}. 
Für softwarebasierte Prozessautomatisierung wird die Auswahl häufig auf Teilaufgaben oder Prozessschritte fokussiert, weil Automatisierung in der Praxis oft nicht die vollständige Ende-zu-Ende-Ausführung ersetzt, sondern wiederkehrende, klar abgrenzbare Tätigkeiten adressiert \cite{AalstBichlerHeinzl2018RPA}. 

Für \gls{rpa} werden Tätigkeiten als besonders geeignet beschrieben, wenn sie regelbasiert, repetitiv und in hoher Frequenz ausgeführt werden \cite{AalstBichlerHeinzl2018RPA}. 
Die Eignung von \gls{rpa} wird zudem mit der Ausführung über bestehende Benutzeroberflächen verknüpft, weil \gls{rpa} in einem „outside-in“-Ansatz menschliche Interaktionen mit vorhandenen Systemen nachbildet \cite{AalstBichlerHeinzl2018RPA}. 
Aus dieser Eigenschaft folgt, dass stabile Oberflächen und geringe Varianz in Eingaben und Prozesspfaden als vorteilhaft eingeordnet werden, weil Änderungen an Oberflächen und häufige Ausnahmebehandlungen den Automatisierungsaufwand erhöhen \cite{Wanner2022ProcessSelectionRPA}. 
Wanner konsolidiert hierfür vier messorientierte Kriterien für das Automatisierungspotenzial von Prozesstätigkeiten, die Ausführungsfrequenz, Stabilität im Sinne geringer Ausnahmen, Regelbasierung der Entscheidungsschritte sowie Standardisierbarkeit der Eingaben und Outputs umfassen \cite{Wanner2022ProcessSelectionRPA}. 
Die Prozessauswahl kann damit als Bewertung verstanden werden, bei der Nutzenpotenziale über Volumen und Bearbeitungsaufwände adressiert werden und Realisierungsrisiken über Varianz, Ausnahmen und Systemabhängigkeiten berücksichtigt werden \cite{Wanner2022ProcessSelectionRPA}. 
Diese Heuristiken sind mit der Positionierung von \gls{rpa} als „lightweight automation“ konsistent, weil der Fokus auf schnell implementierbaren, klar strukturierten Tätigkeiten liegt \cite{AalstBichlerHeinzl2018RPA}. 

Process Mining wird als Ansatz beschrieben, der auf Basis von Ereignisdaten tatsächliche Prozessabläufe sichtbar macht und damit \emph{as-is}-Prozesse empirisch aus Logdaten rekonstruiert \cite{Aalst2016ProcessMining}. 
Van der Aalst ordnet Process Mining dabei nicht nur der Prozessentdeckung zu, sondern auch der Konformitätsprüfung und der Leistungsanalyse, wodurch Abweichungen, Engpässe und Varianten systematisch identifiziert werden können \cite{Aalst2016ProcessMining}. 
Für die Prozessidentifikation unterstützt Process Mining die Auswahl, indem es Häufigkeiten, Variantenvielfalt und Durchlaufzeiten aus Event Logs ableitbar macht und damit Kriterien wie Volumen und Stabilität datenbasiert unterlegt \cite{Aalst2016ProcessMining}. 
Choi et al. beschreiben eine Methodik zur Auswahl von Kandidaten für \gls{rpa}, die Benutzeroberflächen-Logs mit Process-Mining-Techniken kombiniert, um geeignete digitale Aufgaben aus beobachtbarem Ausführungsverhalten abzuleiten \cite{Choi2021CandidateTasksRPA}. 
El-Gharib und Amyot ordnen Process Mining als unterstützenden Baustein im \gls{rpa}-Lebenszyklus ein, weil Process Mining zur Analyse des \emph{as-is}-Prozesses genutzt werden kann und damit die Identifikation automatisierbarer Prozesssegmente strukturieren soll \cite{ElGharibAmyot2023SLR}. 
Die Kombination aus heuristischer Vorauswahl und process-mining-basierter Transparenz adressiert damit die Zielsetzung, Automatisierungskandidaten nachvollziehbar zu priorisieren und zugleich Umsetzungsrisiken über Varianten, Ausnahmen und Datenverfügbarkeit frühzeitig zu erkennen \cite{Wanner2022ProcessSelectionRPA,Aalst2016ProcessMining}. 



\subsection{Governance, Risiken und Compliance}
Governance im Kontext der Prozessautomatisierung beschreibt die Festlegung von Rollen, Verantwortlichkeiten, Entscheidungsrechten und Kontrollmechanismen, um Automatisierungsvorhaben konsistent zu steuern und dauerhaft betreibbar zu halten \cite{Kedziora2021GovernanceModelsRPA}. 
Für \gls{rpa} werden Governance-Modelle unter anderem über die organisatorische Verankerung der Automatisierungsverantwortung und über abgestimmte Regeln für Entwicklung, Betrieb und Änderungsmanagement konkretisiert \cite{Kedziora2021GovernanceModelsRPA}. 
Eine verbreitete organisatorische Ausprägung ist die Bündelung von Kompetenzen und Standards in einem zentralen oder hybriden Steuerungsmodell, das Skalierung und Wiederverwendung von Automatisierungen unterstützt \cite{Kedziora2021GovernanceModelsRPA}. 
Die Governance-Perspektive adressiert damit explizit die Steuerung von Lebenszyklusaktivitäten, die von der Idee über Entwicklung und Test bis zu Betrieb und Stilllegung reichen \cite{Patricio2025AuditFrameworkRPA}. 

Risiken der Prozessautomatisierung werden in der Praxis häufig als Kontroll- und Betriebsrisiken beschrieben, weil automatisierte Ausführungen bestehende Kontrollumgebungen verändern und neue Fehler- sowie Missbrauchsvektoren schaffen können \cite{COSO2025RPAICIF}. 
Für RPA werden interne Kontrollen insbesondere dort als kritisch eingeordnet, wo Automatisierungen in Finanz- und Berichtskontexten eingesetzt werden und damit Anforderungen an Nachvollziehbarkeit und Kontrollwirksamkeit berühren \cite{Kokina2024RPAGovernanceJIS,AAA2024RPAAccountingRiskIC}. 
Risikokategorien betreffen dabei unter anderem unzureichend kontrollierte Berechtigungen von Bots, fehlende Protokollierung, mangelhaftes Change-Management und unzureichende Überwachung der Ausführung, weil diese Faktoren die Integrität von Prozessausgaben und die Revisionsfähigkeit beeinträchtigen können \cite{COSO2025RPAICIF,Kokina2024RPAGovernanceJIS}. 
Ein Audit-orientierter Blick auf RPA ordnet das Risiko zudem entlang des Technologie- und Lebenszykluskontexts ein, weil Kontrollen je nach Phase von Design, Implementierung und Betrieb unterschiedliche Schwerpunkte besitzen \cite{Patricio2025AuditFrameworkRPA}. 

Compliance wird als Fähigkeit verstanden, rechtliche und interne Anforderungen systematisch zu identifizieren, umzusetzen, zu überwachen und fortlaufend zu verbessern \cite{ISO37301}. 
ISO~37301 beschreibt hierfür ein Compliance-Management-System als organisatorischen Rahmen, der die Etablierung, Implementierung, Bewertung, Aufrechterhaltung und Verbesserung von Compliance-Prozessen adressiert \cite{ISO37301}. 
Für Prozessautomatisierung ist Compliance relevant, weil automatisierte Abläufe dieselben Regelwerke und internen Vorgaben erfüllen müssen wie manuell ausgeführte Prozesse und Abweichungen durch Skaleneffekte schneller Wirkung entfalten können \cite{COSO2025RPAICIF}. 
Rechtliche Fragestellungen bei \gls{rpa} werden zudem im Zusammenhang mit Verantwortlichkeit, Transparenz und datenschutzbezogenen Anforderungen diskutiert, weil RPA über Benutzeroberflächen in bestehende Anwendungen eingreift und dabei Daten verarbeitet \cite{Alpers2019LegalChallengesRPA}. 
Ein zentraler europäischer Regulierungsrahmen für die Verarbeitung personenbezogener Daten ist die DSGVO, die Anforderungen an rechtmäßige Verarbeitung und an den Schutz personenbezogener Daten formuliert \cite{GDPR_EURLex}. 

Sicherheits- und Datenschutzanforderungen werden in der Governance-Perspektive häufig mit einem risikobasierten Managementansatz verknüpft, weil Schutzmaßnahmen in Relation zu Bedrohungen und Auswirkungen festgelegt werden \cite{ISO27001}. 
ISO/IEC~27001 beschreibt ein \gls{isms} als Managementsystem, das die Vertraulichkeit, Integrität und Verfügbarkeit von Informationen durch einen Risikomanagementprozess erhält \cite{ISO27001}. 
Für Automatisierungsvorhaben folgt daraus die Notwendigkeit, Bot-Identitäten, Zugriffskonzepte, Protokollierung und Änderungsprozesse so zu gestalten, dass Sicherheitsziele und Compliance-Vorgaben in Betrieb und Weiterentwicklung dauerhaft eingehalten werden können \cite{ISO27001,COSO2025RPAICIF}. 
Governance, Risiken und Compliance bilden damit einen zusammenhängenden Rahmen, weil Governance die organisatorische Steuerung definiert, Risikomanagement die Kontrollschwerpunkte im Lebenszyklus adressiert und Compliance die Einhaltung externer und interner Anforderungen operationalisiert \cite{Kedziora2021GovernanceModelsRPA,ISO37301,COSO2025RPAICIF}. 


\subsection{Herausforderungen und Erfolgsfaktoren in \glspl{kmu}}
Die Einführung von Prozessautomatisierung in \glspl{kmu} ist häufig durch begrenzte personelle und finanzielle Ressourcen geprägt, die den parallelen Aufbau von Fachwissen, Prozessdokumentation und technischer Umsetzung erschweren \cite{OECD2021DT_SME,OECD2021SMEsGoingDigital}. 
Digitale Transformationsvorhaben in \glspl{kmu} werden zudem durch Kompetenzlücken und Weiterbildungsbedarfe beeinflusst, wodurch Implementierungs- und Betriebsaufgaben nicht immer intern abgedeckt werden können \cite{OECD2021DT_SME,OECD2021SMEsGoingDigital}. 
RPA wird als Automatisierungsansatz beschrieben, der über die Benutzeroberfläche bestehender Systeme arbeitet, wodurch Änderungen an Applikationsoberflächen und Prozessvarianten die Stabilität von Automatisierungen beeinträchtigen können \cite{AalstBichlerHeinzl2018RPA,Moreira2023RPAReview}. 
RPA-Projekte adressieren regelbasierte, repetitive Tätigkeiten, sodass hohe Varianz, viele Ausnahmen und unklare Entscheidungsregeln die Auswahl geeigneter Kandidaten und die robuste Umsetzung erschweren \cite{AalstBichlerHeinzl2018RPA,Wanner2019ProcessSelection}. 
Eine ad-hoc- und nachfragegetriebene Umsetzung von RPA wird als wiederkehrende Herausforderung beschrieben, wenn Methoden zur systematischen Identifikation und Priorisierung geeigneter Prozesse fehlen \cite{Wanner2019ProcessSelection}. 
Risiken entstehen zudem, wenn Automatisierungen ohne ausreichende Kontrollen betrieben werden, weil Berechtigungen, Protokollierung und Änderungsmanagement die Integrität der Prozessergebnisse und die Nachvollziehbarkeit der Ausführung beeinflussen \cite{Schlegel2024RatingRisksRPA,COSO2025RPAICIF}. 
Governance-Fragen werden im RPA-Kontext als relevant beschrieben, weil Skalierung, Wiederverwendung und konsistente Entwicklungs- und Betriebsstandards typischerweise organisatorische Regeln und Zuständigkeiten erfordern \cite{Kedziora2021GovernanceModelsRPA,Plattfaut2022CSF_RPA}. 

Erfolgsfaktoren werden zunächst in einer systematischen Prozessauswahl verortet, weil die Eignung von Automatisierungskandidaten von Kriterien wie Regelbasierung, Standardisierbarkeit, Stabilität und ausreichendem Volumen abhängt \cite{Wanner2019ProcessSelection,AalstBichlerHeinzl2018RPA}. 
Ein tragfähiges Zielbild erfordert außerdem Prozessverständnis und Prozessmodellierung, weil BPM Prozesse als Ketten von Ereignissen, Aktivitäten und Entscheidungen beschreibt und Automatisierung eine klare Abgrenzung des betrachteten Prozesssegments voraussetzt \cite{DumasEtAl2013}. 
Plattfaut et al. bündeln kritische Erfolgsfaktoren für RPA in mehreren Kontextclustern, zu denen unter anderem Managementunterstützung, geeignete Governance-Strukturen und kompetenzbezogene Voraussetzungen zählen \cite{Plattfaut2022CSF_RPA}. 
Eine strukturierte Vorgehensweise über den gesamten Lebenszyklus wird als erfolgsrelevant beschrieben, weil RPA-Projekte Faktoren von der Initialisierung über Design und Implementierung bis zum Betrieb und zur kontinuierlichen Verbesserung adressieren müssen \cite{Schlegel2024SuccessfulBPA}. 
Die organisatorische Bündelung von Standards, Methoden und Unterstützungsleistungen wird als Governance-Option diskutiert, weil zentrale oder hybride Modelle die Qualitätssicherung und die Skalierung von Automatisierungen unterstützen sollen \cite{Kedziora2021GovernanceModelsRPA}. 
Kompetenzaufbau und Change-Management werden als erfolgsrelevante Enabler beschrieben, weil fehlende Fähigkeiten und Akzeptanzprobleme die Einführung und nachhaltige Nutzung von Automatisierungslösungen beeinträchtigen können \cite{OECD2021SMEsGoingDigital,Plattfaut2022CSF_RPA}. 
Die Wirksamkeit der Automatisierung wird schließlich durch Betriebskontrollen und Monitoring unterstützt, weil die frühzeitige Erkennung von Ausführungsfehlern und Prozessänderungen die Stabilität automatisierter Abläufe erhöht \cite{Schlegel2024RatingRisksRPA,COSO2025RPAICIF}. 
Für diese Arbeit sind die genannten Herausforderungen und Erfolgsfaktoren relevant, weil sie als Kontextbedingungen für eine ressourcensensitive und risikoangemessene Bewertung von Prozessautomatisierungsinitiativen in \glspl{kmu} dienen.



\section{Synergien zwischen \gls{bi} und Prozessautomatisierung}\label{sec:bi-proc-connection}

\subsection{Integrationsmuster}
Synergien zwischen \gls{bi} und Prozessautomatisierung lassen sich als Integrationsmuster beschreiben, in denen Daten, Analysen und Ausführungslogik wechselseitig gekoppelt werden \cite{DumasEtAl2013,GartnerABIArticle}. 
Ein zentrales Muster ist \emph{BI$\rightarrow$Automation}, bei dem analytische Ergebnisse als Input für die Auswahl, Parametrisierung oder Steuerung von Automatisierungen genutzt werden \cite{GartnerABIArticle,DumasEtAl2013}. 
BI-Plattformen stellen hierfür typischerweise Reporting- und Visualisierungsartefakte wie Dashboards und Berichte bereit, die zur Bewertung von Prozessleistung und zur Ableitung von Handlungsbedarf verwendet werden können \cite{GartnerABIArticle}. 
BPM beschreibt die kontinuierliche Unterstützung des Prozesslebenszyklus einschließlich Analyse, (Re-)Design, Implementierung und Monitoring, wodurch analytische Erkenntnisse als Auslöser für Automatisierungsentscheidungen in diesen Lebenszyklus eingebettet werden können \cite{DumasEtAl2013}. 
Business Activity Monitoring (BAM) wird als Ansatz zur kontinuierlichen, echtzeitnahen Leistungsmessung von Geschäftsprozessen anhand von \glspl{kpi} beschrieben, wodurch Messwerte als Trigger oder Entscheidungsgrundlage für operative Eingriffe genutzt werden können \cite{Wetzstein2008BAM}. 

Ein zweites Muster ist \emph{Automation$\rightarrow$BI}, bei dem Prozessautomatisierung Daten erzeugt, die in BI-Systeme zurückfließen und dort für Monitoring, Ursachenanalyse und Steuerung genutzt werden \cite{Grigori2004BPI,ISO27001}. 
BPMS werden als Plattformen beschrieben, die Prozesse ausführen und dabei Informationen über die unterstützten Prozesse protokollieren, wodurch Ausführungsdaten als Datengrundlage für Analysen verfügbar werden \cite{Grigori2004BPI}. 
Die Analyse von Ausführungslogs wird als Möglichkeit beschrieben, Wissen über Prozessausführung zu gewinnen und damit die Qualität von Prozessen und Services zu verbessern \cite{Grigori2004BPI}. 
RPA wird als „outside-in“-Ansatz beschrieben, bei dem bestehende Informationssysteme unverändert bleiben und Software-Roboter Benutzerinteraktionen nachbilden, wodurch zusätzlich zu Systemlogs häufig Bot-spezifische Laufzeit-, Fehler- und Transaktionsdaten entstehen \cite{AalstBichlerHeinzl2018RPA}. 
Diese Automatisierungsdaten können in BI-Umgebungen zur Transparenz über Automatisierungsleistung, Stabilität und Auslastung genutzt werden, weil BI-Plattformen die Modellierung, Analyse und Visualisierung von Daten zur Entscheidungsunterstützung adressieren \cite{GartnerABIPlatformsReview}. 

Als drittes Muster kann ein \emph{Closed-Loop}-Ansatz beschrieben werden, bei dem Analyse, Ableitung von Maßnahmen und Anpassung der Prozessausführung als Rückkopplungsschleife organisiert werden \cite{Aalst2007FullBPMLifecycle,Wetzstein2008BAM}. 
Process Mining wird als Ansatz beschrieben, der aus Event Logs tatsächliche Prozessabläufe rekonstruiert und hierfür Process Discovery, Conformance Checking und Performance Analysis adressiert \cite{Aalst2016ProcessMining}. 
Diese Fähigkeiten unterstützen Closed-Loop-Mechanismen, weil Abweichungen, Varianten und Engpässe aus Ausführungsdaten identifiziert und als Input für Prozessanpassungen oder Automatisierungsmaßnahmen genutzt werden können \cite{Aalst2016ProcessMining,Aalst2007FullBPMLifecycle}. 
Die Integration von Process Mining und RPA wird als Anwendungsfeld beschrieben, das die Implementierung von RPA beschleunigen und verbessern soll, wodurch datenbasierte Auswahl und kontinuierliche Optimierung als gekoppelte Aktivität umgesetzt werden kann \cite{ElGharibAmyot2023SLR}. 
Closed-Loop-Integrationsmuster erfordern dabei die konsistente Kopplung von Datenbereitstellung, Analyse und Ausführungssteuerung, weil Rückkopplungsschleifen nur dann wirksam sind, wenn Messwerte, Interpretationen und Änderungen in nachvollziehbaren Zyklen zusammengeführt werden \cite{Aalst2007FullBPMLifecycle,DumasEtAl2013}. 


\subsection{Nutzenpotenziale und Zielkonflikte}
Ein Nutzenpotenzial der Kopplung von \gls{bi} und Prozessautomatisierung liegt in einer verbesserten Steuerbarkeit von Prozessausführung, weil Business Process Intelligence explizit Analyse, Monitoring, Vorhersage, Kontrolle und Optimierung der Prozessausführungsqualität adressiert \cite{Grigori2004BPI}. 
Diese Steuerbarkeit wird durch die Nutzung von Ausführungsdaten unterstützt, weil BPMS- und Prozesslogdaten als Grundlage zur Ableitung von Erkenntnissen über tatsächliche Prozessabläufe beschrieben werden \cite{Grigori2004BPI,Aalst2016ProcessMining}. 
Business Activity Monitoring wird als Echtzeitüberwachung von Geschäftsprozessen und -aktivitäten beschrieben, wodurch Kennzahlen als operative Indikatoren für Eingriffe oder Priorisierungen genutzt werden können \cite{IBMBAM}. 
Ein \emph{BI$\rightarrow$Automation}-Muster kann damit Nutzen erzeugen, wenn Kennzahlen und Analysen in den BPM-Lebenszyklus eingebettet werden und Prozessanpassungen sowie Implementierungsentscheidungen datenbasiert ausgelöst werden \cite{DumasEtAl2013,Aalst2007FullBPMLifecycle}. 

Ein weiteres Nutzenpotenzial ergibt sich im \emph{Automation$\rightarrow$BI}-Muster aus zusätzlicher Transparenz über Automatisierungsleistung, weil \gls{rpa} als UI-basierter Ansatz beschrieben wird, der Ausführung, Fehler und Interaktionen softwareseitig abbildet und damit bot-spezifische Laufzeitinformationen generieren kann \cite{AalstBichlerHeinzl2018RPA}. 
Process Mining unterstützt dieses Transparenzpotenzial, weil Prozessentdeckung, Konformitätsprüfung und Performanceanalyse als Kernelemente beschrieben werden, die Abweichungen, Varianten und Engpässe aus Ereignisdaten identifizieren \cite{Aalst2016ProcessMining}. 
Ein \emph{Closed-Loop}-Ansatz kann daraus Nutzen ziehen, weil Prozessanalyse und intelligentes Redesign als zusammenhängende Aktivitäten im BPM-Lebenszyklus beschrieben werden \cite{Aalst2007FullBPMLifecycle}. 

Zielkonflikte entstehen, wenn Automatisierung lokal optimiert wird, während Ende-zu-Ende-Prozessziele unverändert bleiben, weil BPM den Wertbeitrag über zusammenhängende Prozessketten und nicht über isolierte Einzeltätigkeiten beschreibt \cite{DumasEtAl2013}. 
Ein weiterer Zielkonflikt betrifft die Robustheit, weil \gls{rpa} als „outside-in“-Automatisierung beschrieben wird, deren Funktionsfähigkeit von stabilen Oberflächen und konsistenter Interaktion abhängt \cite{AalstBichlerHeinzl2018RPA}. 
Zusätzliche Zielkonflikte betreffen die Interpretierbarkeit von Kennzahlen und die Verlässlichkeit analytischer Steuerungsimpulse, weil BI- und Analytics-Konzepte die Transformation von Daten in entscheidungsrelevante Informationen adressieren und damit die Konsistenz der zugrunde liegenden Daten zur Voraussetzung wird \cite{BaarsKemper2021,WangStrong1996}



\section{Methoden und Technologien für das Zielbild}\label{sec:methods-tech} 

\subsection{Reifegradanalyse für \gls{bi} und Automatisierung}
Reifegradmodelle beschreiben typischerweise mehrere archetypische Stufen, die den Entwicklungsstand einer organisationalen Fähigkeit in einer sequenziellen Logik abbilden und dadurch eine Einordnung des Ist-Zustands unterstützen \cite{RoeglingerPoeppelbuSS2011}. 
Die Nutzung von Reifegradmodellen wird für deskriptive Zwecke als Ist-Assessment und für präskriptive Zwecke als Grundlage zur Ableitung einer Verbesserungsroadmap beschrieben \cite{RoeglingerPoeppelbuSS2011}. 
BPM-bezogene Reifegradmodelle werden als mehrdimensional charakterisiert und umfassen häufig Dimensionen wie Governance, Methoden und Werkzeuge, IT sowie Kultur, wodurch die Bewertung nicht ausschließlich technisch erfolgt \cite{RoeglingerPoeppelbuSS2011}. 
Eine Reifegradanalyse kann damit als strukturierte Methode verstanden werden, die Stärken und Schwächen entlang definierter Dimensionen sichtbar macht und priorisierbare Entwicklungsbedarfe ableitet \cite{RoeglingerPoeppelbuSS2011}. 

Für \gls{bi} existiert eine Vielzahl von Reifegradmodellen, die zur Bewertung von BI-Initiativen herangezogen werden können \cite{LahrmannEtAl2010BIMMsOverview}. 
Lahrmann et al. beschreiben BI-Reifegradmodelle als Instrumente, um Stärken und Schwächen von BI-Bemühungen zu identifizieren und die Auswahl eines geeigneten Modells zu unterstützen \cite{LahrmannEtAl2010BIMMsOverview}. 
Das TDWI-orientierte Modell nach Eckerson wird als stufenbasiertes BI-Reifegradmodell beschrieben, das Reife über mehrere Dimensionen wie Architektur, Daten, Entwicklung und Delivery adressiert \cite{HribarRajteric2010BIMMsOverview,LahrmannEtAl2010BIMMsOverview}. 
Die Stufenlogik wird dabei unter anderem über Metaphern von frühen bis fortgeschrittenen Entwicklungszuständen beschrieben und dient der Einordnung, ob BI eher als fragmentierte Reporting-Lösung oder als unternehmensweite Fähigkeit betrieben wird \cite{HribarRajteric2010BIMMsOverview}. 
Ergänzend wird Analytics-Reife als Stufenmodell beschrieben, das die organisatorische Nutzung von Analytics von einer begrenzten, lokalisierten Nutzung hin zu systematischer, wettbewerbsdifferenzierender Nutzung entwickelt \cite{DavenportHarris2017CompetingOnAnalytics}. 

Für Prozessautomatisierung wird Reife in der Praxis häufig über organisatorische Fähigkeiten wie Prozessauswahl, Entwicklungs- und Betriebsstandards sowie Steuerungs- und Kontrollmechanismen operationalisiert \cite{Anagnoste2013RPACoE,RoeglingerPoeppelbuSS2011}. 
Anagnoste beschreibt im Kontext eines \gls{rpa}-Center-of-Excellence eine Reifedimension, die zwischen Aufbauaspekten wie Prozesspriorisierung, Qualifizierung und Change-Management sowie Betriebsaspekten wie Betriebssicherheit, Change-Handling, Compliance und \glspl{kpi} unterscheidet \cite{Anagnoste2013RPACoE}. 
Für die organisatorische Automatisierungsreife wird zudem ein capability- und levelbasiertes Modell beschrieben, das Reife als Menge von Fähigkeiten und Reifestufen für die Adoption von Best Practices in der Automatisierung adressiert \cite{Microsoft2022AutomationMaturityModel}. 
Eine integrierte Reifegradanalyse für \gls{bi} und Automatisierung kann dadurch entlang kompatibler Dimensionen wie Governance, Daten und IT-Fähigkeiten, Methodenstandardisierung sowie Betrieb und Monitoring strukturiert werden \cite{RoeglingerPoeppelbuSS2011,LahrmannEtAl2010BIMMsOverview,Anagnoste2013RPACoE}. 
Für diese Arbeit unterstützt die Reifegradanalyse die Ableitung eines konsistenten Zielbilds, indem sie Ist-Ausprägungen nachvollziehbar in priorisierte Entwicklungsbedarfe überführt.


\subsection{Entscheidungs- und Bewertungsverfahren}
Entscheidungs- und Bewertungsverfahren für \gls{bi}- und Automatisierungsvorhaben werden häufig als Kombination aus multi-kriterieller Bewertung und wirtschaftlicher Vorteilhaftigkeitsrechnung umgesetzt, weil neben monetären Effekten auch qualitative Zielgrößen wie Transparenz, Steuerbarkeit oder Umsetzbarkeit berücksichtigt werden müssen \cite{BeltonStewart2002MCDA,Zangemeister1976Nutzwertanalyse}. 
Multi-Criteria Decision Analysis (MCDA) beschreibt hierzu Verfahren, die Alternativen anhand mehrerer, potenziell konfligierender Kriterien strukturieren und eine nachvollziehbare Aggregation von Bewertungen ermöglichen \cite{BeltonStewart2002MCDA,KeeneyRaiffa1976DMO}. 
Die \emph{Nutzwertanalyse} wird als Vorgehen zur multidimensionalen Bewertung und Auswahl von Alternativen beschrieben, bei dem Kriterien festgelegt, gewichtet und Alternativen über Skalen bewertet werden, um eine vergleichende Gesamtbewertung zu erzeugen \cite{Zangemeister1976Nutzwertanalyse}. 
Eine solche Nutzwertlogik ist für Technologieentscheidungen geeignet, weil sie qualitative Kriterien explizit macht und damit trade-offs zwischen Zielgrößen transparent abbildet \cite{BeltonStewart2002MCDA,KeeneyRaiffa1976DMO}. 

Ergänzend wird der \emph{Business Case} als Instrument genutzt, um Nutzen, Kosten und Risiken einer Investition strukturiert zu begründen und mit organisatorischen Zielen zu verknüpfen \cite{WardDaniel2006Benefits}. 
Benefits Management beschreibt hierzu einen Prozess, in dem Nutzen identifiziert, geplant, realisiert und anschließend überprüft wird, wodurch Nutzenannahmen in umsetzungsnahe Maßnahmen und Verantwortlichkeiten überführt werden \cite{WardDaniel2006Benefits}. 
Für die monetäre Bewertung wird häufig die Logik des \emph{Net Present Value} (NPV) herangezogen, bei der der Wert eines Projekts über diskontierte zukünftige Cashflows bestimmt und dem Investitionsaufwand gegenübergestellt wird \cite{BrealeyMyersAllenCh5NPV}. 
Die NPV-Regel wird als Entscheidungskriterium beschrieben, nach dem Projekte mit positivem NPV angenommen und Projekte mit negativem NPV abgelehnt werden \cite{BrealeyMyersAllenCh5NPV}. 
Die kombinierte Anwendung von Nutzwertanalyse und Business Case unterstützt damit eine Bewertung, in der qualitative Zielbeiträge strukturiert vergleichbar werden und zugleich die wirtschaftliche Vorteilhaftigkeit über konsistente finanzielle Kriterien abgeleitet wird \cite{Zangemeister1976Nutzwertanalyse,WardDaniel2006Benefits,BrealeyMyersAllenCh5NPV}. 


\subsection{\gls{llm}-gestützte Verarbeitung von Eingaben}
\paragraph{Einordnung und Zielsetzung}
\glspl{llm} basieren in vielen aktuellen Systemen auf der Transformer-Architektur, die Sequenzverarbeitung primär über Attention-Mechanismen realisiert \cite{Vaswani2017}. 
Große, vortrainierte Sprachmodelle werden als skalierte autoregressive Modelle beschrieben, deren Leistungsfähigkeit in instruktionalen und Few-Shot-Szenarien mit zunehmender Modellgröße und Datenmenge steigt \cite{Brown2020}. 
Eine \gls{llm}-gestützte Eingabeverarbeitung kann dadurch als Pipeline verstanden werden, die unstrukturierte oder schwach strukturierte Eingaben in explizite, weiterverarbeitbare Repräsentationen überführt \cite{Brown2020}. 

\paragraph{Eingabetypen und Normalisierung}
Eingaben in BI- und Automatisierungskontexten liegen häufig als Freitext, Formularinhalte, E-Mail-Texte oder Systemmeldungen vor, wodurch zunächst eine Normalisierung der Darstellung und Semantik erforderlich wird \cite{Brown2020}. 
Die Nutzung von \glspl{llm} für diese Normalisierung stützt sich auf deren Fähigkeit, Aufgaben über natürlichsprachliche Instruktionen ohne explizites task-spezifisches Training auszuführen \cite{Brown2020}. 
Für die Weiterverarbeitung wird die Eingabe typischerweise in eine Zielstruktur überführt, damit nachgelagerte Komponenten wie Workflows, Regeln oder Datenmodelle deterministische Schnittstellen erhalten \cite{OpenAIStructuredOutputs}. 

\paragraph{Informationsextraktion und Strukturierung}
Die Extraktion von Entitäten, Attributen und Relationen aus Texten kann als Kernleistung der \gls{llm}-basierten Eingabeverarbeitung eingeordnet werden, weil natürliche Sprache in strukturierte Felder transformiert werden soll \cite{OpenAIStructuredOutputs}. 
Strukturierte Ausgaben können über eine vorgegebene JSON-Schema-Spezifikation erzwungen werden, wodurch Schlüssel, Datentypen und zulässige Wertebereiche explizit begrenzt werden \cite{OpenAIStructuredOutputs}. 
Die Schema-Vorgabe reduziert Integrationsaufwand, weil nachgelagerte Systeme nicht auf freie Textausgaben reagieren müssen, sondern valide, maschinenlesbare Objekte erhalten \cite{OpenAIStructuredOutputs}. 
Die technische Kopplung an operative Komponenten kann zusätzlich über Tool- bzw. Function-Calling realisiert werden, bei dem das Modell einen strukturierten Aufruf samt Parametern erzeugt, der anschließend durch eine externe Funktion ausgeführt wird \cite{OpenAIFunctionCalling}. 

\paragraph{Kontextanreicherung durch Retrieval-Augmented Generation}
Vortrainierte Sprachmodelle speichern Wissen in Parametern, weisen jedoch Grenzen bei Aktualisierbarkeit, Präzision der Wissensnutzung und der Bereitstellung von Provenienz auf \cite{Lewis2020RAG}. 
Retrieval-Augmented Generation (RAG) kombiniert parametriertes Modellwissen mit nichtparametrischem Wissen aus einem externen Index, der zur Laufzeit abgefragt wird \cite{Lewis2020RAG}. 
RAG wird als Architektur beschrieben, in der ein Retriever relevante Passagen selektiert und ein Generator die Ausgabe unter Konditionierung auf diese Passagen erzeugt \cite{Lewis2020RAG}. 
Für die Eingabeverarbeitung in BI- und Automatisierungsszenarien unterstützt RAG die Domänenanpassung, weil organisationsspezifische Dokumente, Richtlinien oder Datenbeschreibungen zur Laufzeit als Kontext eingebunden werden können \cite{Lewis2020RAG,LangChainRetrieval}. 
Die praktische Umsetzung von Retrieval erfordert üblicherweise eine Segmentierung großer Dokumente in kleinere Einheiten, damit relevante Kontexte in das begrenzte Modellkontextfenster passen \cite{LangChainSplitters}. 

\paragraph{Closed-Loop-Kopplung in der Eingabeverarbeitung}
Ein Closed-Loop-Ansatz kann auf Eingabeebene dadurch entstehen, dass aus Extraktionen und Klassifikationen unmittelbar Aktionen in Workflows ausgelöst werden und Ausführungsdaten anschließend zur Verbesserung der Eingabemodelle und Regeln zurückfließen \cite{OpenAIFunctionCalling,Lewis2020RAG}. 
Die Kopplung von strukturierter Extraktion und Tool-Ausführung unterstützt diesen Loop, weil Eingaben in maschinenlesbare Entscheidungen überführt und anschließend kontrolliert ausgeführt werden können \cite{OpenAIStructuredOutputs,OpenAIFunctionCalling}. 
RAG kann den Loop zusätzlich stabilisieren, weil Referenztexte als Kontext die Konsistenz der Interpretation erhöhen und Änderungen an Wissensbeständen ohne erneutes Modelltraining abbildbar werden \cite{Lewis2020RAG,LangChainRetrieval}. 

\paragraph{Risiken, Robustheit und Governance-Aspekte}
Prompt-Injection wird als Sicherheitsrisiko für \gls{llm}-integrierte Anwendungen beschrieben, bei dem Angreifer Anweisungen in Eingaben verstecken, um das Modellverhalten zu manipulieren \cite{Liu2023PromptInjection}. 
Dieses Risiko ist besonders relevant, wenn Eingaben aus externen Kanälen stammen und das Modell Zugriff auf Tools oder Systeme erhält, weil manipulative Instruktionen sonst in ungewollte Aktionen übersetzt werden können \cite{Liu2023PromptInjection,OpenAIFunctionCalling}. 
Risikomanagement für generative KI wird im NIST AI RMF als organisationsübergreifende Aufgabe beschrieben, die Kontext, Auswirkungen und geeignete Kontrollen systematisch berücksichtigen soll \cite{NISTAIRMF1001}. 
Der NIST-Generative-AI-Profile-Rahmen ordnet generative KI explizit als Anwendungsfeld des AI RMF ein und beschreibt die Notwendigkeit, Risiken entlang des gesamten Lebenszyklus zu adressieren \cite{NISTGenAIProfile6001}. 
Für die Eingabeverarbeitung folgt daraus die Notwendigkeit, Output-Validierung, Zugriffskontrollen und Protokollierung als technische Leitplanken zu etablieren, wenn Modelloutputs operative Prozesse steuern \cite{NISTAIRMF1001,OpenAIStructuredOutputs}. 



\subsection{Process \& Data Mining als Enabler}
Process Mining wird als Ansatz beschrieben, der aus Ereignisdaten (\emph{event logs}) tatsächliche Prozessabläufe rekonstruiert und damit eine empirische Sicht auf die Prozessausführung ermöglicht \cite{Aalst2016ProcessMining}. 
Die Kernaufgaben von Process Mining umfassen Process Discovery, Conformance Checking und Performance Analysis, wodurch Varianten, Abweichungen und Engpässe systematisch identifizierbar werden \cite{Aalst2016ProcessMining}. 
Diese Transparenz unterstützt die Auswahl von Automatisierungskandidaten, weil Häufigkeiten, Durchlaufzeiten und Ausnahmequoten aus Logs als objektivierbare Indikatoren für Standardisierbarkeit und Stabilität genutzt werden können \cite{Aalst2016ProcessMining,AalstBichlerHeinzl2018RPA}. 
Process Mining kann zudem als Teil von Business Process Intelligence eingeordnet werden, weil Business Process Intelligence Analyse, Monitoring und Optimierung von Prozessausführung auf Basis von Ausführungsdaten adressiert \cite{Grigori2004BPI}. 

Data Mining wird als Teil des Knowledge Discovery in Databases (KDD) beschrieben, bei dem aus Datenbeständen Muster und Modelle abgeleitet werden, um aus Daten verwertbares Wissen zu erzeugen \cite{Fayyad1996KDD}. 
Der KDD-Prozess wird als Sequenz von Schritten wie Selektion, Vorverarbeitung, Transformation, Data Mining und Interpretation beschrieben, wodurch Datenaufbereitung und Modellbildung als zusammenhängende Pipeline operationalisiert werden \cite{Fayyad1996KDD}. 
Data-Mining-Verfahren umfassen unter anderem Klassifikation, Clustering und Assoziationsanalyse, die in Standardwerken als zentrale Aufgabenklassen der Datenanalyse dargestellt werden \cite{HanKamberPei2012}. 
In BI- und Automatisierungskontexten kann Data Mining als Enabler wirken, weil Mustererkennung und Prognosemodelle zur Priorisierung von Fällen, zur Erkennung von Anomalien oder zur Ableitung von Handlungsempfehlungen genutzt werden können \cite{HanKamberPei2012,Grigori2004BPI}. 

Die kombinierte Nutzung von Process Mining und Data Mining unterstützt Closed-Loop-Ansätze, weil Logdaten sowohl zur Rekonstruktion des Prozessverhaltens als auch zur Ableitung von prädiktiven oder erklärenden Modellen herangezogen werden können \cite{Aalst2016ProcessMining,HanKamberPei2012}. 
Für die technische Umsetzung sind dabei konsistente Ereignisprotokollierung und Datenqualität als Voraussetzung zu verstehen, weil die Güte der abgeleiteten Modelle von der Verlässlichkeit der zugrunde liegenden Daten abhängt \cite{Aalst2016ProcessMining,WangStrong1996}. 


\subsection{Monitoring, Definition von \glspl{kpi} und Erfolgsmessung}
Monitoring wird im BPM-Kontext als fortlaufende Beobachtung der Prozessausführung und der Prozessleistung verstanden, die auf Messwerten aus der Ausführung basiert und damit die Steuerung entlang des Prozesslebenszyklus unterstützt \cite{DumasEtAl2013}. 
Eine solche Beobachtung stützt sich auf die systematische Erhebung und Auswertung von Leistungsinformationen, wie sie im Rahmen von Business Process Intelligence für Analyse, Monitoring und Optimierung der Prozessausführung beschrieben wird \cite{Grigori2004BPI}. 
Business Activity Monitoring adressiert hierfür die kontinuierliche, potenziell echtzeitnahe Ermittlung von Prozesskennzahlen, wodurch Messwerte zeitnah zur Verfügung stehen \cite{Wetzstein2008BAM}. 

\glspl{kpi} werden als Kennzahlen verstanden, die auf kritische Aspekte der organisatorischen Leistung fokussieren und dadurch Steuerungsimpulse ermöglichen \cite{Parmenter2015KPI}. 
Die Definition von \glspl{kpi} wird als Aufgabe beschrieben, bei der Kennzahlen aus strategischen Zielsetzungen abgeleitet werden, damit Messgrößen den intendierten Zielbeitrag abbilden \cite{KaplanNorton1996BSC}. 
Die Balanced Scorecard ordnet Kennzahlen entlang mehrerer Perspektiven, um strategische Ziele in messbare Größen zu übersetzen und dadurch Steuerung zu ermöglichen \cite{KaplanNorton1996BSC}. 
Für die operative Anwendbarkeit von \glspl{kpi} ist eine eindeutige Spezifikation erforderlich, weil Messdefinitionen ohne klare Berechnungslogik, Datenquelle, Aggregationsregeln und Verantwortlichkeiten nicht zuverlässig vergleichbar sind \cite{Parmenter2015KPI}. 
Die modellbasierte Erfassung von \glspl{kpi} in BAM-Ansätzen wird als Vorgehen beschrieben, bei dem \glspl{kpi} so definiert werden, dass sie in eventbasierten Ausführungsmodellen messbar sind \cite{Wetzstein2008BAM}. 

Erfolgsmessung wird in IS- und IT-Investitionskontexten als Bestandteil eines Benefits-Management-Prozesses beschrieben, der die Identifikation, Planung, Realisierung und Überprüfung von Nutzen adressiert \cite{WardDaniel2006Benefits}. 
Die Überprüfung von Nutzen erfordert dafür definierte Messgrößen und eine Zuordnung von Verantwortlichkeiten, damit Nutzenannahmen in überprüfbare Ergebnisse überführt werden können \cite{WardDaniel2006Benefits}. 
Im Zusammenspiel von \gls{bi} und Automatisierung unterstützt die Kopplung von Monitoring und \glspl{kpi} die Bewertung von Integrationsentscheidungen, weil Prozess- und Automatisierungsdaten als Grundlage für Steuerung und Priorisierung beschrieben werden \cite{Grigori2004BPI,Wetzstein2008BAM}. 
Für die Zielsetzung dieser Arbeit sind Monitoring, \gls{kpi}-Definition und Erfolgsmessung relevant, weil sie Kriterien und Messmechanismen bereitstellen, um Alternativen nachvollziehbar zu bewerten und Maßnahmen in einem Entscheidungsframework konsistent zu operationalisieren.


\section*{Zusammenfassung}
Optional?